\documentclass[12pt]{article}
\usepackage[utf8]{inputenc}
\usepackage[spanish]{babel}
\usepackage[a4paper, margin=2cm]{geometry}
\usepackage{tikz}
\usetikzlibrary{babel}
\usepackage[most]{tcolorbox}
\usepackage{amsmath}

% Usar fuente sin serifa para un aspecto más moderno (tipo infografía)
\renewcommand{\familydefault}{\sfdefault}

% Definición de colores con temática médica/científica
\definecolor{medblue}{RGB}{43, 108, 176}
\definecolor{medteal}{RGB}{49, 151, 149}
\definecolor{medlight}{RGB}{235, 248, 255}
\definecolor{meddark}{RGB}{26, 32, 44}
\definecolor{petriblue}{RGB}{191, 219, 254}   % Azul suave para la placa de Petri
\definecolor{bacteriagreen}{RGB}{16, 185, 129} % Verde para las bacterias
\definecolor{zoomborder}{RGB}{156, 163, 175}   % Gris para las líneas de zoom

\begin{document}
\pagestyle{empty}

\begin{tcolorbox}[
    enhanced, 
    colback=medlight!30, 
    colframe=medblue, 
    title={\Large \textbf{Potencias de Diez: Biomasa Celular}}, 
    halign title=center, 
    fonttitle=\bfseries, 
    arc=4mm, 
    boxrule=1.5pt,
    shadow={2mm}{-2mm}{0mm}{black!30!white}
]

    \vspace{0.2cm}
    \begin{tcolorbox}[colback=white, colframe=medteal, arc=2mm, boxrule=1.2pt, title=\textbf{Caso Clínico / Problema}]
        En un cultivo de laboratorio, se cuenta una colonia de bacterias con \textbf{$2 \times 10^6$} individuos. Si se asume que cada bacteria tiene una masa de \textbf{$1 \times 10^{-12}$} gramos, ¿cuál es la masa total de la colonia?
    \end{tcolorbox}

    \vspace{0.4cm}
    
    \begin{center}
    \begin{tikzpicture}[scale=0.95]
        
        % Placa de Petri (Izquierda)
        \filldraw[petriblue!40, draw=medblue, thick] (-3,0) circle (1.5cm);
        \node[font=\bfseries, text=meddark, align=center] at (-3,-2.2) {Colonia de Bacterias\\$2 \times 10^6$ individuos};
        
        % Dibujar pequeñas bacterias en la placa de Petri
        \foreach \a in {0, 45, 90, 135, 180, 225, 270, 315} {
            \filldraw[bacteriagreen, draw=meddark, very thin, rotate around={\a:(-3,0)}] (-2.5,0) ellipse (0.15cm and 0.08cm);
            \filldraw[bacteriagreen, draw=meddark, very thin, rotate around={\a+20:(-3,0)}] (-2.0,0) ellipse (0.15cm and 0.08cm);
        }
        \filldraw[bacteriagreen, draw=meddark, very thin, rotate around={10:(-3,0)}] (-2.8,0) ellipse (0.15cm and 0.08cm);
        
        % Líneas de efecto Zoom
        \draw[dashed, zoomborder, thick] (-2.7, 0.2) -- (1.5, 1.2);
        \draw[dashed, zoomborder, thick] (-2.7, -0.2) -- (1.5, -1.2);
        
        % Círculo de Zoom (Derecha)
        \filldraw[white, draw=zoomborder, line width=2pt] (3,0) circle (1.5cm);
        
        % Bacteria individual en el zoom
        \filldraw[bacteriagreen, draw=meddark, line width=1.5pt, rounded corners=5pt] (2.0, -0.5) rectangle (4.0, 0.5);
        % Flagelos (colitas de la bacteria)
        \draw[meddark, line width=1.5pt] (2.0, 0.1) to[out=180,in=90] (1.3, -0.3);
        \draw[meddark, line width=1.5pt] (2.0, -0.1) to[out=180,in=-90] (1.1, 0.4);
        \draw[meddark, line width=1.5pt] (2.0, 0.3) to[out=180,in=0] (1.5, 0.5);
        
        \node[font=\bfseries, text=meddark, align=center] at (3,-2.2) {Masa Individual\\$1 \times 10^{-12}$ g};
        
        % Ecuación visual en la parte inferior
        \node[font=\Large, text=meddark] at (0, -4.5) {
            $\underbrace{(2 \times 10^6)}_{\text{Cantidad}} \times \underbrace{(1 \times 10^{-12})}_{\text{Masa Individual}} \quad \xrightarrow{\text{Agrupar}} \quad \underbrace{(2 \times 1) \times (10^6 \times 10^{-12})}_{\text{Coeficientes y Bases}}$
        };
        
    \end{tikzpicture}
    \end{center}

    \vspace{0.4cm}
    \begin{tcolorbox}[colback=medteal!10, colframe=medteal, arc=2mm, boxrule=1.2pt, title=\textbf{Desarrollo Matemático y Solución}]
        Para encontrar la masa total de la biomasa, multiplicamos la cantidad de bacterias por la masa de una de ellas. Agrupamos los coeficientes y las bases:
        \begin{align*}
            M &= (2 \times 1) \times 10^{6 + (-12)} \\[1ex]
            M &= 2 \times 10^{-6} \text{ g}
        \end{align*}
        \textbf{Resultado:} La masa total de la colonia es de \textbf{$2 \times 10^{-6}$ g}, lo que equivale en el Sistema Internacional a \textbf{2 microgramos ($\mu$g)}.
    \end{tcolorbox}
    
\end{tcolorbox}

\end{document}